% TOP_PROBLEM: {"id": 30006715, "title": "Graph Isomorphism in Polynomial Time", "category_id": 15, "category": {"id": 15, "name": "computer_science", "display_name": "Computer Science", "description": "Computational complexity, algorithms, and theoretical CS.", "slug": "computer-science", "order_index": 15, "created_at": "2026-07-31T15:26:25.671Z"}, "status": "open"}
% FIRST_POSTED: 2026-09-27
% !TeX program = lualatex
% Compile twice: lualatex 30006715.tex
% Standalone research notebook; original section preserved.
% Research additions: 27 September 2026. No external TeX input or BibTeX file.
\documentclass[11pt,a4paper]{article}
\usepackage[margin=24mm,headheight=14pt]{geometry}
\usepackage{fontspec}
\setmainfont{Latin Modern Roman}
\setsansfont{Latin Modern Sans}
\setmonofont{Latin Modern Mono}
\usepackage{amsmath,amssymb,mathtools}
\usepackage{unicode-math}
\setmathfont{Latin Modern Math}
\usepackage{microtype}
\usepackage{enumitem,xurl,needspace}
\usepackage[dvipsnames]{xcolor}
\usepackage{hyperref}
\hypersetup{unicode=true,colorlinks=true,linkcolor=MidnightBlue,urlcolor=MidnightBlue,
 pdftitle={Research notebook - Problem 30006715},pdfauthor={Open Problems Math},bookmarksnumbered=true}
\usepackage{bookmark}
\usepackage{titlesec}
\titleformat{\section}{\Large\bfseries\raggedright}{\thesection}{0.7em}{}
\usepackage{fancyhdr}
\pagestyle{fancy}\fancyhf{}
\fancyhead[L]{\small\sffamily Open Problems Math | Research notebook}
\fancyhead[R]{\small\sffamily Problem 30006715 | 27 September 2026}
\fancyfoot[C]{\thepage}
\setlength{\parindent}{0pt}
\setlength{\parskip}{5pt plus 1pt}
\setlength{\emergencystretch}{3em}
\setcounter{tocdepth}{1}
\setcounter{secnumdepth}{1}
\setlist[itemize]{leftmargin=3em,itemsep=5pt,topsep=4pt,parsep=0pt}
\allowdisplaybreaks
\numberwithin{equation}{section}
\newcommand{\N}{\mathbb{N}}
\newcommand{\Z}{\mathbb{Z}}
\newcommand{\Q}{\mathbb{Q}}
\newcommand{\R}{\mathbb{R}}
\newcommand{\C}{\mathbb{C}}
\newcommand{\F}{\mathbb{F}}
\newcommand{\A}{\mathbb{A}}
\newcommand{\PP}{\mathbb P}
\newcommand{\E}{\mathbb{E}}
\newcommand{\eps}{\varepsilon}
\newcommand{\dd}{\,\mathrm d}
\newcommand{\sref}[2]{\hyperlink{source:#1:#2}{[S#2]}}
\newcommand{\eref}[2]{\hyperlink{extra:#1:#2}{[E#2]}}
\newif\ifshowcatalogue
\showcataloguetrue
\newcommand{\cataloguescope}[1]{\ifshowcatalogue
 \paragraph{Exact scope in the supplied catalogue.}
 \begin{quote}\small #1\end{quote}\fi}
\begin{document}
\setcounter{section}{0}
\section*{\texorpdfstring{Graph Isomorphism in Polynomial Time}{Graph Isomorphism in Polynomial Time}}\label{30006715}
\subsection{Definitions and mathematical statement}
Two finite graphs $G=(V,E)$ and $H=(W,F)$ are isomorphic if some bijection $\varphi:V\to W$ satisfies
\[
 \{u,v\}\in E\quad\Longleftrightarrow\quad\{\varphi(u),\varphi(v)\}\in F.
\]
The problem asks for a deterministic classical algorithm deciding the existence of such a bijection in time polynomial in the ordinary graph encoding length. No planarity, degree bound, or structural restriction is assumed. A quasipolynomial algorithm, or an algorithm efficient only on typical graphs, is not the demanded bound.
\par\smallskip\textit{Formulation and concept references:} \sref{30006715}{1}.
\cataloguescope{Determine whether the graph isomorphism problem for general finite graphs is decidable by a deterministic classical Turing machine in polynomial time.}
\subsection{Short English statement}
Can arbitrary finite graphs be tested for isomorphism in deterministic polynomial time?
% BEGIN RESEARCH_ADDITION ATTEMPT_79
\subsection{Research attempt}

\paragraph{Current frontier.}
Babai's general algorithm is quasipolynomial, not polynomial in the ordinary
encoding length \eref{30006715}{1}. Polynomial or efficient smoothed-case and
numerical methods do not by themselves resolve worst-case deterministic
GI; the smoothed setting in \eref{30006715}{2} illustrates the distinction.
The general polynomial-time claim retained as \sref{30006715}{4} is not treated
as established merely because it was posted. The following attempt tests
a concrete convex-relaxation route and gives an exact stopping certificate
for a sufficiently integral solution.

\paragraph{Attack route.}
Spectral invariants are cheap but incomplete. Fractional isomorphism adds
entrywise positivity and linear constraints, yet permits many nonpermutation
solutions. A plausible repair would be to drive a feasible matrix near a
permutation using a provable integrality gap. The attempt proves a precise
rounding threshold, then constructs two explicit nonisomorphic graphs whose
fractional and heat-kernel information defeats the simpler route. The gap
certificate is mathematical; an efficient method for attaining it is not
assumed.

\paragraph{Attempt: a robust integrality certificate.}
Let $A,B$ be the $n\times n$ adjacency matrices of two simple graphs.
Consider the rational polytope
\[
 \mathcal P(A,B)=\{X\geq0:X\mathbf1=X^{\mathsf T}\mathbf1=\mathbf1,
                                     \ AX=XB\}.
\]
A permutation matrix in this set is an exact isomorphism. Define
$\Delta(X)=n-\|X\|_F^2$. Each row is a probability vector, so
$\Delta(X)\geq0$, with equality exactly when every row is a coordinate
vector. The column sums then make those coordinates distinct.

There is also a nonzero exact threshold:
\begin{equation}\label{eq79:round}
 X\in\mathcal P(A,B),\qquad \Delta(X)<\frac1{4n^2}
 \quad\Longrightarrow\quad G\cong H.
\end{equation}
For a proof, put $m_i=\max_jX_{ij}$. Since
$\sum_jX_{ij}^2\leq m_i$, a row with $m_i\leq1/2$ would already
contribute at least $1/2$ to $\Delta$, impossible under the displayed
bound. The unique entries exceeding $1/2$ in different rows must occupy
distinct columns, by the column-sum condition. They define a permutation
matrix $P$. Moreover,
\[
 \sum_i(1-m_i)\leq\Delta(X),\qquad
 \|X-P\|_F^2=2\sum_i(1-m_i)-\Delta(X)\leq\Delta(X).
\]
Since $\|A\|_{\rm op},\|B\|_{\rm op}\leq n$ and $AX=XB$,
\[
 \|AP-PB\|_F\leq2n\|P-X\|_F<1.
\]
But $AP-PB$ has integer entries. Any nonzero integer matrix has Frobenius
norm at least one, so $AP=PB$. This proves \eqref{eq79:round}.
It yields the contrapositive gap: for a nonisomorphic pair every feasible
$X$ has $\Delta(X)\geq1/(4n^2)$.

The exactness of the linear constraints matters. A rationally certified
feasible matrix below the threshold is a finite, polynomially checkable
isomorphism certificate and can be rounded directly. A floating-point
residual is not automatically such a certificate; quantitative residual
terms would need to be included in the last inequality. Conversely,
searching for the minimum of the concave function $n-\|X\|_F^2$ over the
polytope is not linear programming. The existence of a gap does not prove
that a deterministic polynomial-time algorithm finds a point attaining
it whenever a permutation exists.

\paragraph{Attempt: explicit fractional and spectral obstruction.}
Every pair of $d$-regular graphs of the same order admits
$X=J/n\in\mathcal P(A,B)$, since $AJ=dJ=JB$. Thus feasibility alone
cannot settle GI. For a stronger stress test take two graphs on $16$
vertices. The rook graph has vertices $(i,j)\in(\Z/4\Z)^2$, with edges
between distinct vertices in the same row or column. The Shrikhande graph
is the Cayley graph on the same group with connection set
\[
 \{(1,0),(-1,0),(0,1),(0,-1),(1,1),(-1,-1)\}.
\]
Both are six-regular, and direct neighbour counting gives
\begin{equation}\label{eq79:srg}
 A^2=4I+2J
\end{equation}
for each adjacency matrix: the diagonal is six and each pair of distinct
vertices has two common neighbours. On the subspace orthogonal to
$\mathbf1$, \eqref{eq79:srg} gives eigenvalues $2$ and $-2$; trace zero
then gives multiplicities $6$ and $9$, with the remaining eigenvalue $6$.
So their adjacency spectra coincide exactly.

They are not isomorphic. A row of the rook graph is a four-clique. In the
Shrikhande graph the induced graph on the six neighbours of $(0,0)$ is a
six-cycle, as one checks by subtracting pairs in the displayed connection
set. A clique containing $(0,0)$ therefore has at most three vertices.
Translations are automorphisms, so this holds for a clique through any
vertex; triangles exist. Its clique number is three, proving the claimed
nonisomorphism. The companion script verifies the adjacency identities,
the six-cycle and the clique numbers by exact finite enumeration.

For the graph Laplacian $L=6I-A$, every analytic function of $L$ lies in
$\operatorname{span}\{I,A,J\}$ by \eqref{eq79:srg}. In particular the
heat matrix $e^{-tL}$ has, for every fixed $t$, entries determined solely
by equality, adjacency, or nonadjacency of the two vertices, with identical
three values for both graphs. Hence their global entry multisets agree at
every heat time. This does not make the matrices conjugate by a permutation;
it shows that the multiset invariant throws away the combinatorial
placement of the entries. Also $AX=XB$ automatically implies
$A^rX=XB^r$ for every $r$, so appending those redundant linear constraints
does not remove $J/n$.

\paragraph{Stress test.}
The gap \eqref{eq79:round} could be useful if one had a polynomial-time
global near-integrality method, but no such method is supplied. Local
improvement of a quadratic objective or convergence of a numerical flow
to a stationary point is not a worst-case global certificate. Nor is an
exponential branch search made polynomial simply because most branches
collapse on typical graphs. Those are algorithmic missing steps, not
statements that optimization methods cannot ever help.

The rook/Shrikhande example defeats the specified relaxation and spectral
summaries; it does not defeat all higher-dimensional refinement algorithms,
all invariant rings, or Babai's group-theoretic method. Likewise it is not
a complexity lower bound. An isomorphism algorithm is allowed to exploit
information beyond eigenvalues, heat-entry multisets, and fractional
intertwiners. For graphs of different orders rejection is immediate; the
argument's equal-order assumption is not a restriction on the full decision
problem because that case is handled first.

\paragraph{Outcome.}
\textbf{Outcome: Exploratory approach with unresolved gap.}

A robust exact rounding criterion is proved, and an explicit nonisomorphic
pair defeats a natural convex/spectral simplification. The remaining
question is whether a deterministic polynomial-time procedure can find
near-integral feasible intertwiners whenever an isomorphism exists, or
use a different complete invariant with polynomial evaluation cost.
No such procedure or superpolynomial lower bound is obtained. These are
method tests and a certificate lemma, not a claim of a new general GI
algorithm or of literature priority.
% END RESEARCH_ADDITION ATTEMPT_79
\subsection{Sources}
\begin{itemize}\raggedright
\item[\textnormal{[S1]}]\hypertarget{source:30006715:1}{} E. W. Weisstein, 'Graph Isomorphism,' MathWorld, accessed 2026.
\url{https://mathworld.wolfram.com/GraphIsomorphism.html}.
{\small\textit{Catalogue formulation source.}}
\item[\textnormal{[S2]}]\hypertarget{source:30006715:2}{} Graph Isomorphism: Mixed-Integer Convex Optimization from First-Order Methods.
\url{https://arxiv.org/abs/2512.17417}.
{\small\textit{Catalogue research/status reference; not a proof endorsement.}}
\item[\textnormal{[S3]}]\hypertarget{source:30006715:3}{} Finding Graph Isomorphisms in Heated Spaces in Almost No Time.
\url{https://arxiv.org/abs/2601.03787}.
{\small\textit{Catalogue research/status reference; not a proof endorsement.}}
\item[\textnormal{[S4]}]\hypertarget{source:30006715:4}{} Testing Isomorphism of Graphs in Polynomial Time.
\url{https://arxiv.org/abs/2305.12688}.
{\small\textit{Catalogue research/status reference; not a proof endorsement.}}
% BEGIN RESEARCH_ADDITION REFERENCES_79
\item[\textnormal{[E1]}]\hypertarget{extra:30006715:1}{} L\'aszl\'o Babai,
\emph{Graph Isomorphism in Quasipolynomial Time}, arXiv:1512.03547,
main algorithmic theorem and complexity convention.
\url{https://arxiv.org/abs/1512.03547}.
{\small\textit{Inspected 27 September 2026; quasipolynomial is not polynomial.}}
\item[\textnormal{[E2]}]\hypertarget{extra:30006715:2}{} Michael Anastos,
Matthew Kwan and Benjamin Moore, \emph{Smoothed analysis for graph isomorphism},
STOC 2025, DOI 10.1145/3717823.3718173.
\url{https://research-explorer.ista.ac.at/record/20007}.
{\small\textit{Institutional publication record inspected 27 September
2026; smoothed hypotheses are not dropped in discussing the general target.}}
% END RESEARCH_ADDITION REFERENCES_79
\end{itemize}

\end{document}
